% concatenated from sn-article.tex
\documentclass{amsart}

\usepackage{graphicx}
\usepackage{multirow}
\usepackage{amsmath,amssymb,amsfonts}
\usepackage{mathrsfs}
\usepackage[title]{appendix}
\usepackage{xcolor}
\usepackage{textcomp}
\usepackage{manyfoot}
\usepackage{booktabs}
\usepackage{algorithm}
\usepackage{algorithmicx}
\usepackage{algpseudocode}
\usepackage{listings,hyperref}

\newtheorem{theorem}{Theorem}[section]
\newtheorem{proposition}[theorem]{Proposition}
\newtheorem{lemma}[theorem]{Lemma}
\newtheorem{corollary}[theorem]{Corollary}
\newtheorem{remark}[theorem]{Remark}
\newtheorem{example}[theorem]{Example}
\newtheorem{definition}[theorem]{Definition}
\newtheorem*{theorema}{Theorem}
\newcommand{\NN}{\mathbb N}
\newcommand{\ZZ}{\mathbb Z}
\newcommand{\RR}{\mathbb R}
\newcommand{\CC}{\mathbb C}

\title[Gap estimates for the spectrum of the $m$-bonacci numbers]{Gap estimates for the spectrum of $m$-bonacci numbers}


\author{Anna Chiara Lai}
\address{
Dipartimento di Scienze di Base e Applicate per l'Ingegneria,
Sapienza Universit\`a di Roma,
Via Antonio Scarpa 16,
00161 Roma, Italy}
\email{annachiara.lai@uniroma1.it}

\author{Paola Loreti}
\address{
Dipartimento di Scienze di Base e Applicate per l'Ingegneria,
Sapienza Universit\`a di Roma,
Via Antonio Scarpa 16,
00161 Roma, Italy}
\email{paola.loreti@uniroma1.it}

\subjclass[2020]{11K16, 11A63}

\keywords{Pisot number, golden number, $m$-bonacci number, spectrum, gap conditions}

\begin{document}

\begin{abstract}
We establish some results on the structure of the spectrum of $m$-bonacci numbers. More precisely, we study explicit lower bounds for the distance between elements separated by $N$ positions in the ordered spectrum. The result is further detailed in the Fibonacci and Tribonacci case. The methodology combines the combinatorial structure of $m$-bonacci words and the canonical $m$-bonacci number system.
\end{abstract}

\maketitle

\section{Introduction}
% The analysis of the spectrum of Pisot numbers has been studied for a long time because of its applications to number theory, symbolic dynamics and harmonic analysis.  
In this paper we are interested in gap estimates in the spectrum of $m$-bonacci numbers. 

To illustrate our results, we begin by considering the spectrum $\Lambda (\varphi)$ of the golden mean $\varphi$, namely the set of real numbers of the form $\sum_{i=0}^k \varepsilon_i \varphi^i$
with $k\in \mathbb{N}$ and $\varepsilon_i \in \{0,1\}$. Ordering $\Lambda(\varphi)$ strictly increasingly, so that $\Lambda(\varphi)=\{0<\lambda_1(\varphi)<\cdots<\lambda_n(\varphi)<\cdots\}$ we get
$$\Lambda(\varphi)=\{0,1,\varphi,\varphi^2,\varphi^2+1,\varphi^2+\varphi,\varphi^2+\varphi+1,\dots\}$$ and the consecutive distances (or gaps) $\lambda_{n+1}(\varphi)-\lambda_{n}(\varphi)$ take only two possible values $1$ and $\varphi-1$. If one considers larger increments $\lambda_{n+2}(\varphi)-\lambda_{n}(\varphi)$, the possible values are $2$ and $\varphi$. In particular, there is no $n\in\NN$ such that $\lambda_{n+2}(\varphi)-\lambda_{n}(\varphi)=2\varphi-2$. This constraint can be clearly understood by coding the gap sequence $\lambda_{n+1}(\varphi)-\lambda_n(\varphi)$ by assigning the symbol $1$ or $2$ depending on whether the gap is equal to $1$ or $\varphi-1$. The resulting sequence is $121121\cdots$, which coincides with the Fibonacci word. The Fibonacci word can be defined as the limit of the substitution $1 \mapsto 12$, $2 \mapsto 1$, obtained by iterating this rule starting from the symbol $1$ -- see \cite{AS03}.  In particular, the subword $22$ never occurs; this excludes the possibility of having two consecutive gaps equal to $\varphi-1$, and consequently $\lambda_{n+2}(\varphi)-\lambda_n(\varphi) \neq 2\varphi - 2$ for all $n$. More generally, finer combinatorial properties of the Fibonacci word,
such as balance and forbidden patterns, yield gap estimates of the form
\begin{equation}\label{lambdaest}\lambda_{n+N}(\varphi)-\lambda_{n}(\varphi)\geq N\left(\frac{1}{\varphi}+\frac{1}{\varphi^3}\right) -\left(\frac{1}{\varphi^4}+\varphi\right)\end{equation}
for all $N\geq 2,\,n\geq 1$ -- see Corollary \ref{cor1}. Writing the lower bound in \eqref{lambdaest} in the form $N\gamma_{2,N}$ with $\gamma_{2,N}:=\left(\frac{1}{\varphi}+\frac{1}{\varphi^3}\right) -\frac{1}{N}\left(\frac{1}{\varphi^4}+\varphi\right)$, we obtain that $\gamma_{2,N}\to \frac{1}{\varphi}+\frac{1}{\varphi^3}=\frac{\sqrt{5}}{\varphi}=\frac{1}{D^+(\Lambda(\varphi))}$ as $N\to \infty$, where $D^+(\Lambda(\varphi))$ denotes the (upper) density of $\Lambda(\varphi)$  -- see the definitions below. 

Our goal is to establish estimates of the form \eqref{lambdaest} for the wider class of $m$-bonacci numbers $q_m$, defined as the largest roots of the polynomials 
$$x^m-x^{m-1}-\cdots -x-1 \qquad \text{for } m\geq 2.$$
Note that the case $m=2$ corresponds to the golden mean case, because $q_2=\varphi$. We consider the spectrum of the $m$-bonacci numbers
$$\Lambda(q_m):=\left\{\sum_{i=0}^n\varepsilon_iq_m^i\mid \varepsilon_{i}\in\{0,1\},\, n=0,1,\dots\right\}=\{0<\lambda_1(q_m)<\lambda_2(q_m)<\cdots\}$$
and we establish gap estimates  of the form 
$$\lambda_{n+N}(q_m)-\lambda_{n}(q_m)\geq N\gamma_{m,N} \quad \text{for all }n\geq 1$$
 for any $N\geq 1$. Our main result provides a constant $\gamma_{m,N}$ which can be interpreted as a weighted average of the gap values in $\Lambda(q_m)$. The weights depend on the $m$-bonacci expansion of $N$,  and on the combinatorial structure of
the $m$-bonacci word coding the gaps. 
An $m$-bonacci expansion is a representation
of $N$ as a sum 
 of $m$-bonacci numbers $F^{(m)}_n$ with coefficients in $\{0,1\}$. The $m$-bonacci numbers are recursively defined by 
$F^{(m)}_0=1$ $F^{(m)}_1=2, \cdots F^{(m)}_{m-1}=2^{m-1}$
and, for $n\geq m$,
$$F^{(m)}_n=F^{(m)}_{n-1}+\cdots +F^{(m)}_{n-m}.$$
On the other hand, we recall that the $m$-bonacci word is defined as the fixed point of the Rauzy
substitution $\sigma_m$ defined by 
$$1\mapsto 12,\quad 2\mapsto 13,\quad \dots, \quad m-1\mapsto 1m, \quad m\mapsto 1.$$
In particular, the weights involve the balance constant $b_m$, which measures how uniformly the letters are distributed the $m$-bonacci word. 
Our main result is the following.  
\begin{theorem}\label{thm1i}
Let $m\geq 2$ and $N\in \mathbb{N}$. Then there exists a constant
$\gamma_{m,N}>0$ such that
\begin{equation}\label{gamma}
\lambda_{n+N}(q_m)-\lambda_n(q_m)\geq N\,\gamma_{m,N}
\quad \text{for all } n\geq 1.
\end{equation}

The following choice of $\gamma_{m,N}$ satisfies the above inequality.
Let $(x_h)\in\{0,1\}^\infty$ be the canonical $m$-bonacci expansion of $N$, that is,
a binary sequence with no occurrences of the block $1^m$ that expands $N$ in the form
\[
N=\sum_{h\geq 0} x_h F^{(m)}_h,
\quad x_h\in\{0,1\},
\]
where $(F^{(m)}_h)$ denotes the $m$-bonacci sequence.

Then
\begin{equation}\label{eqgamma}
\gamma_{m,N}
=\frac{1}{N}
\sum_{j=1}^m
\left(
\sum_{h\geq 0} x_h F^{(m)}_{h-j}
- b_m
\right)
d_m(j),
\end{equation}
where:
\begin{itemize}
\item[-] $d_m(j):=q_m^{j-1}-\sum_{k=0}^{j-2} q_m^k$ with $j=1,\dots,m$ are the finitely many possible gap values in $\Lambda(q_m)$;
\item[-] $b_m$ is a constant lower or equal than $m-1$ measuring the balance of the $m$-bonacci word.
\end{itemize}
\end{theorem}
We remark that explicit values for $b_m$ can be assured by results on the balance of $m$-bonacci word, borrowing from \cite{bal,balancetrib,generalbalance} we have that \eqref{eqgamma} holds by setting 
\begin{equation}\label{bm}
b_m:=\begin{cases}
    m-1 & \text{for } m=2,3,4;\\
    \lceil \frac{m+1}{2}\rceil& \text{for } m=5,\dots,12;\\
    \lfloor \kappa m\rfloor+12 & \text{for } m>12
\end{cases}\end{equation}
where $\kappa:= \frac{2}{\pi} \int_{0}^{2\pi}
\frac{1-\cos x}{(5-4\cos x)\ln(5-4\cos x)} \, dx \approx 0.58.$



By applying the above result to $m=2$ and remarking that the balance constant in the Fibonacci word is $b_2=1$, we obtain in the golden mean case 
$$\lambda_{N+n}(\varphi)-\lambda_{n}(\varphi)\geq \gamma_{2,N}=
\frac{1}{N}
\left(
\sum_{h=0}^\infty x_h(\varphi F_{h-1}+F_{h-2})
-\varphi^2
\right).$$
where $(x_h)\in\{0,1\}^\infty$ is the canonical Fibonacci expansion of $N$. Note that, in general, $m$-bonacci expansions $(x_h)$ 
have all $x_h=0$ except for finitely many. In Corollary \ref{cor1} and in Corollary \ref{cor2} we discuss in detail the cases $m=2,3$, respectively. 
\medskip


% \emph{State of the art.}
% In non-harmonic analysis, expressions of the form $\lambda_{n+N}(q)-\lambda_{n}(q)\geq N\gamma_N$ (in the above inequality $\gamma_N=1-\varphi/N$) are weakened gap conditions strongly related to Ingham-type inequalities. 
Let us place our contribution in the context of the literature on the spectra associated with Pisot numbers and gap distribution. 
The paper \cite{EJK90} is among the first to investigate the spectrum of a real number $q>1$, defined as:
$$\Lambda(q):=\left\{\sum_{i=0}^n\varepsilon_iq^i\mid \varepsilon_{i}\in\{0,1\},\, n=0,1,\dots\right\}=\{0<\lambda_1(q)<\lambda_2(q)<\cdots\}.$$
 Erd{\"o}s, Jo{\'o}, and Komornik were interested in the gap size $\lambda_{n+1}(q)-\lambda_n(q)$ showing that it is uniformly bounded from above by $1$ for all $q\in (1,2)$. The interest then moved to lower estimates: in the same paper it is proved that if $\lambda_{n+1}(q)-\lambda_n(q)\to 0$   then $1$ has an infinite expansion containing
arbitrarily long sequences of consecutive $0$ digits, namely $1=\sum_{i=1}^\infty q^{n_i}$ for some diverging sequence $n_i$ of natural numbers satisfying $\limsup_{i\to\infty} (n_{i+1}-n_i)=+\infty$. 
These results paved the way to many other investigations. In \cite{EJ92}, Erd{\"o}s, Jo{\'o}, and Jo{\'o}  showed that the $m$-bonacci number  $q_m$  is the unique root of $x^m-x^{m-1}-\cdots-x-1$ in $[1,2]$, $q_m\to2$ monotone increasingly as $m\to \infty$. They considered 
the finite sets
$\Lambda^{(h)}(q):=\left\{\sum_{k=1}^h \varepsilon_iq^i\mid \varepsilon_i\in\{0,1\}\right\}$
  for $h\geq 1$ and proved that the minimum gap in $\Lambda^{(h)}(q_m)$ is greater or equal to $\frac{1}{q_m}$ and equality holds for $h\geq m+1$.  As remarked in \cite{EJK98}, this result can be interpreted as follows. Setting for all $q>0$
  $$\ell(q):=\liminf_{n} \lambda_{n+1}(q)-\lambda_{n}(q)$$
  we obtain 
  $$\ell(q_m)=\frac{1}{q_m} \quad \text{for all }m\geq 2. $$
  In \cite{bug96}, Bugeaud considers the wider spectrum \begin{equation}\label{genspectrum}\Lambda^r(q):=\left\{\sum_{i=0}^n\varepsilon_iq^i\mid \varepsilon_{i}\in\{0,1,\dots,r\},\, n=0,1,\dots\right\}=\{0<\lambda_1^r(q)<\cdots<\lambda_n^r(q)\}\end{equation}
and the related quantity 
$$\ell^r(q):=\liminf_{n} (\lambda^r_{n+1}(q)-\lambda^r_n(q)).$$
Bugeaud shows that a real number $q\in(1,2)$ is a \emph{Pisot number}, namely a real algebraic integer  whose algebraic conjugates have modulus stricly less than $1$, if and only if $\ell^r(q)>0$ for all $r\geq 1$. Note that, for $m\geq 2$, $q_m$ is a Pisot number so Bugeaud's result applies. In \cite{EJK98} it is  proved that if $q>1$ is Pisot, then $\ell(q)>0$.
In the paper \cite{KLP00}, Komornik, Loreti, and Pedicini provide a complete description of $\ell^r(\varphi)$ when $\varphi$ is the golden mean: if $F_k$ is the $k$-th Fibonacci number $F_0=0,F_1=1,\, F_{k}=F_{k-1}+F_{k-2}$, and $\varphi^{k-2}< r\leq \varphi^{k-1}$ then $\ell^r(\varphi)=|F_{k}\varphi-F_{k+1}|$. 
% More recently, Feng  shows in \cite{F16} that for all $r\geq 1$ we have that $\ell^r(q)>0$ if and only if either $q$ is Pisot or $r<q-1$.


Now, in \cite{FW02}, Feng and Wen investigate the structure of the set $\Lambda(q)$ in the Pisot case, showing that the gap sequence $\lambda_{n+1}(q)-\lambda_n(q)$ admits a finite number of values and that it can be coded via substitutions on a finite alphabet. A substitution on a finite alphabet assigns to each symbol a finite word over the same alphabet, and is then extended to longer words by concatenation. As mentioned above, by \cite{FW02,bug02}, 
when $q=q_m$ is the largest root of $x^m-x^{m-1}-\cdots-x-1$, the gaps in $\Lambda(q_m)$ are related to the Rauzy substitution $\sigma_m$
via its fixed point 
$$\mathbf v_{m}=\mathbf v_{m,1}\mathbf v_{m,2}\cdots:=\lim_{n\to\infty} \sigma_m^n(1)$$
that encodes the combinatorial structure of the spectrum. Explicit gap formulae are given in \cite{bug02} and, as corollary, Bugeaud proved in the same paper that the gap $1$ occurs in $\Lambda(q_m)$ with frequency $\frac{1}{q_m}$ whereas the gap 
$q_m^{v-1}-q_m^{v-2}-\cdots q_m-1$ occurs with frequency $\frac{1}{q_m^{v}}$. In \cite{H04} Hare gives an algorithm for the computation of the frequencies of gap sizes in more general cases -- this result is generalized to a wider classes of Pisot numbers in the paper \cite{GH06} by Garth and Hare. Akyiama and Komornik \cite{AK13} prove that the set generalized spectrum $Y^r(q):=\{\sum_{i=0}^n\varepsilon_i q^i\mid \varepsilon_i\in\{0,\pm 1,\dots,\pm r\}\}$ has not accumulation points if and only if $q$ is a Pisot number or $q\geq m+1$.   More recently, Feng  shows in \cite{F16} that for all $r\geq 1$ we have that $\ell^r(q)>0$ if and only if either $q$ is Pisot or $r<q-1$. 

Finally, we mention that the spectrum $\Lambda(q)$ of a Pisot number
is closely related to several fundamental objects in aperiodic order
and numeration theory. To illustrate this connection, we note that
$\Lambda(q_m)$ and the $m$-bonacci quasicrystal can be described in terms of the same substitution, although they give rise to different sets of gap sizes. For an overview, we refer to the survey by Berthé and Yassawi \cite{BY24} and the references therein.
\medskip

\emph{Organization of the paper.}
In Section \ref{sec2} we collect the preliminary material needed for the proof,
with particular emphasis on the combinatorial structure of the
$m$-bonacci word, and gap frequencies.
In Section \ref{sec3} we prove Theorem~\ref{thm1i}, combining the
$m$-bonacci expansions of integers with the balance properties of the $m$-bonacci word,  and we derive further estimates  for the Fibonacci case and the Tribonacci case. 
In Section \ref{sec4} we explore some density properties of
$\Lambda(q_m)$. 
Finally, in Section \ref{sec5} we present our conclusions.
\section{Preliminaries}\label{sec2}
We recall the definitions and results on $m$-bonacci words, $m$-bonacci expansions and density of $\Lambda(q_m)$.

\subsection{The $m$-bonacci words and their relation with $\Lambda(q_m)$}
For any $m\geq 2$, the $m$-bonacci word $\mathbf v_m$ is generated iteratively, via the repeated application of a substitution map $\sigma_m$ --known as the \emph{Rauzy substitution}-- up to convergence to its fixed point, which is $\mathbf v_m$. 

\begin{definition}[Rauzy substitution and $m$-bonacci word]   Let $\mathcal A_m:=\{1,\dots,m\}$ with $m\geq 2$ and $\mathcal A_m^*$ be the set of finite words with digits in $\mathcal A_m$. For every $m\geq 2$, we consider \emph{Rauzy substitution} $\sigma_m:\mathcal A_m^*\to \mathcal A_m^*$ given by
$k\mapsto 1(k+1)$ for $k=1,\dots,m-1$ and $m\mapsto 1$. 


For all $m\geq 2$ let
    $$\mathbf v_m:=\lim_{n\to \infty}\sigma^n_m(1).$$ 
    
   The words $\mathbf v_2=\sigma_2^\infty(1)=121121\cdots$  and $\mathbf v_3=\sigma_3^\infty(1)=121312\cdots$ are respectively called the \emph{Fibonacci word} and \emph{Tribonacci word}. For $m> 3$, the word $\mathbf v_m$ is called the \emph{$m$-bonacci word}. 
\end{definition}
For a general introduction on substitutions and, in particular, on the well-posedness of the limit in the definition of $\mathbf v_m$ we refer to \cite{Loth}. Here we limit ourselves to remarking that the substitution $\sigma_m(1)$ begins with $1$. 
Hence the sequence of words $(\sigma_m^n(1))_{n\geq 1}$ is increasing
for the prefix order and thus converges to a unique infinite word.

\medskip
The next result relates $m$-bonacci words to the gaps in the spectrum $\Lambda(q_m)$ in terms of size and frequency. Recall that a gap $g$ occurs with frequency $f$ in an increasingly ordered set $\{\lambda_k\}$ if 
$$\lim_{n\to \infty}\frac{\#\{k\mid \lambda_{k+1}-\lambda_{k}=g, \, k\leq n\}}{n}=f.$$


\begin{theorem}[Bugeaud \cite{bug02}]\label{thm2}
Let $m \geq 2$ and let $q_m$ be the largest real root in $(1,2)$ of
\[
x^m - x^{m-1} - \cdots - x - 1.
\]
Let $\mathbf v_m = \mathbf v_{m,1}\mathbf v_{m,2}\cdots$ be the
$m$-bonacci word. 
Then, for all $n \geq 1$, one has
\[
\lambda_{n+1}(q_m)-\lambda_n(q_m)=d_m(j)
\quad \text{if and only if} \quad
\mathbf v_{m,n}=j,
\]
where
\[
d_m(1)=1,
\qquad
d_m(j)=q_m^{j-1}-q_m^{j-2}-\cdots - q_m - 1
\quad \text{for } j=2,\dots,m.
\]

Moreover, for $j=1,\dots,m$ the value $d_m(j)$ occurs in the gap sequence
$\bigl(\lambda_{n+1}(q_m)-\lambda_n(q_m)\bigr)_{n\geq 1}$
with frequency
$q_m^{-j}.$
\end{theorem}
For instance, in the golden mean case $m=2$, the gaps $1$ and $\varphi-1$ occur in $\{\lambda_{n+1}(\varphi)-\lambda_{n+1}(\varphi)\}$ with frequencies $1/\varphi$ and $1/\varphi^2$, respectively. 

\begin{example}[Fibonacci word]
The Fibonacci word corresponds to the case $m=2$. The first iterations of $\sigma_2$ on $1$ are
\begin{align*}
&\sigma_2(1)=12;\\
&\sigma_2^2(1)=\sigma_2(12)=\sigma_2(1)\sigma_2(2)=121;;\\
&\sigma_2^3(1)=12112;\\
&\sigma_2^4(1)=12112121\end{align*}
and for all $k$, $\sigma_2^k(1)$ is a prefix of $\mathbf v_2$. 
% The PF eigenvalue $\rho_2$ is the Golden Ratio $\varphi:=(1+\sqrt{5})/2$, namely the greatest root of $x^2-x-1$. 
We have $d_2(1)=1$ and $d_2(2)=\varphi-1=1/\varphi$. 
The first three positive terms  of the gap sequence are 
$$\lambda_1(\varphi)-\lambda_0(\varphi)=d_2(\mathbf v_{2,1})=d_2(1)=1$$
and
$$\lambda_2(\varphi)-\lambda_1(\varphi)=d_2(\mathbf v_{2,2})=d_2(2)=\frac{1}{\varphi},\qquad \lambda_3(\varphi)-\lambda_2(\varphi)=d_2(\mathbf v_{2,3})=d_2(1)=1$$ Then 
$$\Lambda(\varphi)=\{0,1,1+1/\varphi,2+1/\varphi,\dots\}.$$


%Incidentally, note that the length of $\sigma_2^n(1)$ is the $n+1$-th Fibonacci number and this confirms that $\sigma_2^\infty(1)=\mathbf w_2$ is infinite.

% Since the normalized positive left eigenvalue of the adjacency matrix $M_{\sigma_2}$ is the vector $(\varphi^{-1},\varphi^{-2})$ ---because the Golden mean is the Perron-Frobenius eigenvalue of $M_{\sigma_2}$-- then the frequencies of the digits $1$ and $2$ in $\mathbf w_2$ are respectively $\varphi^{-1}$ and $\varphi^{-2}$.  
\end{example}

\begin{example}[Tribonacci word]
The first iterations of $\sigma_3$ on $1$ are
\begin{align*}
&\sigma_3(1)=12;\\
&\sigma_3^2(1)=\sigma_3(12)=\sigma_3(1)\sigma_3(2)=1213;\\
&\sigma_3^3(1)=1213121;\\
&\sigma_3^4(1)=1213121121312. \end{align*}
and for all $k$, $\sigma_3^k(1)$ is a prefix of $\mathbf v_3$. We have $\rho_3=\tau$, where $\tau$ is the greatest root of $x^3-x^2-x-1$, i.e., the Tribonacci constant. We have
$$d_3(1)=1\,\quad d_3(2)= \tau-1,   \quad d_3(3)=\tau^2-\tau-1.$$


The first three positive terms of the gap sequence are
\[
\lambda_1(\tau)-\lambda_0(\tau)
=
d_3(\mathbf v_{3,1})
=
d_3(1)
=
1,
\]
and
\[
\lambda_2(\tau)-\lambda_1(\tau)
=
d_3(\mathbf v_{3,2})
=
d_3(2)
=
\tau-1,
\qquad
\lambda_3(\tau)-\lambda_2(\tau)
=
d_3(\mathbf v_{3,3})
=
d_3(1)
=
1.
\]

Therefore
\[
\Lambda(\tau)
=
\{0,1,\tau,\tau+1,\dots\}.
\]
\end{example}


\subsection{Combinatorics on the $m$-bonacci word}
Theorem \ref{thm2} relates the occurrences of the gaps $d_m(j)$ in $\Lambda(q_m)$ to the digit distibution in the $m$-bonacci word. Therefore, it is useful in the sequel to state some definitions and results on the combinatorial properties of $\mathbf v_m$.

\begin{definition}[Weight and balanced word]
Let $w$ be a finite or infinite word over the alphabet $\mathcal A_m=\{1,\dots,m\}$.
For $j\in \mathcal A_m$, we denote by $|w|_j$ the number of occurrences of the letter $j$ in $w$, and call it the \emph{weight} of $j$ in $w$.

Let $b$ be a positive integer. An infinite word $\mathbf v$ over $\mathcal A_m$ is said to be $b$-\emph{balanced} for any two subwords $u$ and $w$ of $\mathbf v$ of the same length, and for every $j\in \mathcal A_m$, one has
\[
\bigl||u|_j - |w|_j\bigr| \leq b.
\]
\end{definition}
The next result describes the weights of the $m$-bonacci word. 
\begin{lemma}\label{l1bal}[\cite{balancetrib}]
For every prefix $v$ of the $m$-bonacci word $\mathbf v_m$ there exist $H$ and $x_0,x_1,\dots,x_H\in\{0,1\}$ such that 
\begin{equation}\label{prefix}(|v|_1,\dots,|v|_m)^T=\sum_{h=0}^H x_h M_m^h (1,0,\dots,0)^T \end{equation}
where $M_m\in \RR^m\times \RR^m$ is the incidence matrix of the Rauzy substitution $\sigma_m$, namely
\begin{equation}\label{Mm}M_{m}= 
\begin{pmatrix}
1 & 1 & 1 & \cdots & 1 & 1 \\
1 & 0 & 0 & \cdots & 0 & 0 \\
0 & 1 & 0 & \cdots & 0 & 0 \\
\vdots & \vdots & \vdots & \ddots & \vdots & \vdots \\
0 & 0 & 0 & \cdots & 1 & 0
\end{pmatrix}.
\end{equation}
Moreover, for every choiche of $H$ and $x_0,\cdots,x_H\in\{0,1\}$ there exists a prefix of $\mathbf v_m$ such that \eqref{prefix} holds. 
\end{lemma}

The $m$-bonacci words are balanced, the next result collects some estimates on the balance constants $b_m$ -- note that the case $m\geq 3$ makes use of the above result.

\begin{theorem}[\cite{bal,balancetrib,generalbalance}]\label{thm3}
For all $m\geq 2$ the $m$-bonacci word is $b$-balanced with 
\begin{enumerate}
\item $b=m-1$ for $m=2,3,4$ and this bound cannot be improved.
\item $b=\lceil \frac{m+1}{2}\rceil$ for $m=5,\dots,12$ and if $m=5$ this bound cannot be improved.
\end{enumerate}
Morevoer for all $m\geq 5$ the following estimate hold
$$b=\lceil \kappa m\rceil+12$$ 
where $\kappa:= \frac{2}{\pi} \int_{0}^{2\pi}
\frac{1-\cos x}{(5-4\cos x)\ln(5-4\cos x)} \, dx \approx 0.58.$
\end{theorem}
%  In the statement of Theorem \ref{thm1} we set  
% \begin{equation}\label{bm}
% b_m:=\begin{cases}
%     m-1 & \text{for } m=2,3,4;\\
%     \lceil \frac{m+1}{2}\rceil& \text{for } m=5,\dots,12;\\
%     \lfloor \kappa m\rfloor+12 & \text{for } m>12
% \end{cases}\end{equation}
% as a balance constant.
% Remark that in general $b_m\leq m-1$.


\subsection{The $m$-bonacci expansions}
Define the \emph{$m$-bonacci sequence} recursively
$$F^{(m)}_0=1, F^{(m)}_1=2, \cdots F^{(m)}_{m-1}=2^{m-1}$$
and for $n\geq m$
$$F^{(m)}_n=F^{(m)}_{n-1}+\cdots +F^{(m)}_{n-m}.$$
For $m=2$ one obtains the sequence $(F^{(2)})_n$ is $ 1,2,3,5,\dots$, the classical Fibonacci sequence (number A000045  in the OEIS classification) shifted by two indices. 
Similarly,  for $m=3$ the sequence $(F^{(3)}_n)= 1, 2, 4, 7,,\dots$ corresponds to the Tribonacci numbers (number A000073 in the OEIS classification), shifted by three indices.

In \cite{CSH72}, investigate the $m$-bonacci numeration system, also called the higher order Fibonacci expansions. Generalizing earlier results established in the Fibonacci case, the Zeckendorf representation \cite{Lek,zec}, a result in that paper states that every $N\in\NN$ admits an expansion of the form 
\begin{equation}\label{Nfib}N=\sum_{h=0}^{\infty}x_h F^{(m)}_h\qquad x_h\in\{0,1\}\end{equation}
where only finitely many digits $x_h$ are non zero and the digit sequence $(x_h)$ contains no blocks of $m$ consecutive $1$'s. Precisely, the greedy algorithm is the following.


Let $N\in\mathbb N$ and set $R_0:=N$.
For each $t\geq 0$, as long as $R_t>0$, define
\[
h_t:=\max\{h\geq 0:\ F_h^{(m)}\leq R_t\},
\]
and
\[
R_{t+1}:=R_t-F_{h_t}^{(m)}.
\]

Since $R_t$ is strictly decreasing, there exists $T\geq 0$ such that
$R_{T+1}=0$. The greedy $m$-bonacci expansion of $N$ is then
\[
N=\sum_{t=0}^{T}F_{h_t}^{(m)}
   =\sum_{h\geq 0}x_hF_h^{(m)},
\]
where
\[
x_h=
\begin{cases}
1,& \text{if } h=h_t \text{ for some } t,\\[2mm]
0,& \text{otherwise.}
\end{cases}
\]





% Such digit expansion is called \emph{(higher order) canonical expansion}, and it takes also the name of  \emph{$m$-bonacci greedy expansion}. 

Canonical representations are not unique in general. For instance, when $m=2$ we have 
$$3=F_0^{(2)}+F_1^{(2)}=F_2^{(2)}$$
so that $11(0)^\infty$ and $001(0)^\infty$ are both expansions of $3$. 



% we mean  the  lexicographic greatest (or greedy) expansion of $N$. 

 
%  To this end, we need to exploit the combinatorics and the finer structure of $\Lambda(q_m)$. Indeed, our argument is based on
%   exploiting the $m$-bonacci representation of the gap size $N$ in an original way, and on the following preliminary facts: the gaps in $\Lambda(q_m)$ are finite and coded by a balanced, infinite word. More precisely our starting point is that, as remarked in \cite{FW02} and formally proved in \cite{bug02}, 
%  the gaps in $\Lambda(q_m)$ are coded by the \emph{Rauzy substitution} $\sigma_m$ defined by
% $$1\mapsto 12,\quad 2\mapsto 13,\quad \dots, \quad m-1\mapsto 1m, \quad m\mapsto 1$$
% in the following sense. 
% Consider the fixed point 
% $$\mathbf v_{m}=\mathbf v_{m,1}\mathbf v_{m,2}\cdots:=\lim_{n\to\infty} \sigma_m^n(1),$$
% and note that the cases $m=2,3$ yield the Fibonacci and Tribonacci word, respectively. In order to state our main result properly, we recall that $\mathbf v_m$ is a \emph{balanced word}, see \cite{bal,balancetrib,generalbalance} and \cite{Loth} for a general introduction, namely there exists a constant $b_m\leq m-1$ measuring how far two factors of the same length can differ in their letter frequencies: this property will control fluctuations in the distribution of gap types. 


% In  \cite{bug02} Bugeaud showed that the gaps in $\Lambda(q_m)$ can be written as
% \[
% \lambda_{n+1}(q_m)-\lambda_n(q_m) = d_m(j) \quad \text{if and only if } \mathbf v_{m,n}=j\]
% where 
% $$d_m(j):=\begin{cases}
% 1 \quad & \text{if }j=1\\
% q_m^{j-1}-q_m^{j-2}-\cdots q_m-1 \quad &\text{if  }  j\geq 2.
% \end{cases}$$

% Define the \emph{$m$-bonacci sequence} recursively
% $$F^{(m)}_0=1, F^{(m)}_1=2, \cdots F^{(m)}_{m-1}=2^{m-1}$$
% and for $n\geq m$
% $$F^{(m)}_n=F^{(m)}_{n-1}+\cdots +F^{(m)}_{n-m}.$$
% For $m=2$ the sequence $(F^{(2)})_n$ is $ 1,2,3,5,\dots$, namely a shift of two indexes of the the classical Fibonacci sequence (number A000045  in the OEIS classification)  whereas for $m=3$ the sequence $(F^{(3)}_n)= 1, 2, 4, 7,,\dots$ corresponding to a shift of three indexes of the Tribonacci numbers (number A000073 in the OEIS classification).

% Recall from \cite{mnacci} that every $N\in\NN$ admits an expansion of the form 
% $$N=\sum_{h=0}^{\infty}x_h F^{(m)}_h\qquad x_h\in\{0,1\}$$
% with all $x_h=0$ except for finitely many.  In general this expansion is not unique, so we we refer to $(x_h)$ as  \emph{a} $m$-bonacci expansion.
% % we mean  the  lexicographic greatest (or greedy) expansion of $N$. 

% % The next result relates the first digits of $\mathbf v_m$ to the $m$-nacci expansions. 

% \begin{theorem}
% %[Gap conditions ]
% \label{thm1}
% Let $m\geq 2$, $N\in \NN$. For any choice $(x_h)$ of a $m$-bonacci expansion of $N$ the constant 
% $$\gamma_{m,N}:=\frac{1}{N}\left(\sum_{j=1}^m\left(\ \sum_{h=0}^\infty (x_h F^{(m)}_{h-j})-b_m\right)d_m(j)\right)$$
% satisfies the gap condition
% \begin{equation}\label{gengaptr}
% \lambda_{m,N+k}-\lambda_{m,k}\geq N\gamma_{m,N} \quad \text{for all } k\in \ZZ.
% \end{equation}
% \end{theorem}
% By applying the above result to $m=2$ and remarking that the balance constant in the Fibonacci word is $b_2=1$, we obtain in the golden mean case 
% $$\lambda_{N+n}(\varphi)-\lambda_{n}(\varphi)\geq \gamma_{2,N}=
% \frac{1}{N}
% \left(
% \sum_{h=0}^\infty x_h(\varphi F_{h-1}+F_{h-2})
% -\varphi^2
% \right).$$
% \medskip



% The density of the spectrum $\Lambda(q)$ is closed related to results in non-harmonic analysis in the following sense. 


% The interest in the gap condition \eqref{e1} in a control theoretic perspective is mainly motivated by the results in \cite{bkl02} (see also \cite{komlorbook} and the references therein) showing the equivalence of  Beurling's theorem \cite{beu} with a weakened gap condition of the form \eqref{e1}. More precisely, let us assume \eqref{e1} and fix arbitrarily $\gamma'\in(0,\gamma_M)$. Define the following index sets
% \begin{align*}A_j:=\{ &m\in\ZZ\mid \lambda_m-\lambda_{m-1}\geq \gamma',\\
% &\lambda_n-\lambda_{n-1}<\gamma' \text{ for }n\in[m+1,m+j-1]\\
% &\lambda_{m+j}-\lambda_{m+j-1}\geq \gamma'
% \}.
% \end{align*}
% for $j=1,\dots,M$. Note that the collection $\{A_j+k\mid k\in\ZZ, j=1,\dots,M\}$ partitions $\ZZ$.
% For $m\in A_j$ and $n=m+1,\dots,m+j-1$ define the domain 
% \begin{align*}\Omega_n
% :=\{(&s_1,\dots,s_{n-m})\in \RR^{n-m}\mid \\
% &s_1\in[0,1], s_h\in[0,s_{h-1}], h=2,\dots,n-m\}\end{align*}
% and the functions
% \begin{align*}
%     e_n(t):=&(it)^{n-m}\cdot\\
%     &\int_{\Omega_n} 
%     \exp(i\sum_{h=1}^{n-m}s_h(\lambda_{h}-\lambda_{h-1})t)d s_{n-m}\cdots d s_1. 
% \end{align*}
% Finally set
% $$f(t):=\sum_{n\in\ZZ} a_n e_n(t).$$
% with complex coefficients $a_n$. Then by \cite[Theorem 1.5]{bkl02}
% for every interval $I$ of length $|I|>2\pi/ \gamma_M$ there exists two positive constants 
% $c_1,c_2$ such that
% \begin{equation}\label{ingh}c_1\sum_{n\in\ZZ}|a_n|^2\leq \int_{|I|} |f(t)|^2dt \leq c_2 \sum_{n\in \ZZ} |a_n|^2.\end{equation}
% The constants depend only on $M,\gamma_M$ and on the interval $I$. By a classical result by Beurling \cite{beu}  there exists a positive constant $c_1$ satisfying the inequality 
% \begin{equation}
% \int_{I} |\sum_{n}a_n e^{i\lambda_n(q_m) t}|^2dt \leq c_2 \sum_{n\in \ZZ} |a_n|^2.\end{equation}




% When $(\lambda_k)$ is strictly increasing, then $f(t)$ coincides with the trigonometric sum $\sum_{k}a_n e^{i\lambda_n t}$, and the inequalities \eqref{ingh}
% were earlier established for a wide class of quasicrystals (including the $\Lambda_m$) in \cite{matmey,meyer}. Moreover, since the density of $\Lambda_m$ is explicitly derivable from the spectral theory of substitutions \cite{queffel, substitutionberthebook}, analogous results can be  by applying Beurling's theorem \cite{beu}. 

% there

% By a result of 






 

% \textcolor{blue}{Considerare anche le densità del corollario 4.2? aggiungere qui il teorema il beurling?}





% Indeed, in the Pisot case, and in particular for $q=q_m$, this density can be
% expressed in terms of the frequencies of the letters in the associated
% substitutive sequence, which are given by the Perron–Frobenius eigenvector
% of the Rauzy substitution (see, e.g., \cite{queffel,substitutionberthebook,baakegrimm}).





% is .  Classical examples include the golden mean and, more generally, 
% . 

% Moreover the finer structure of the sequence $(\lambda_n)$ has also been clarified by Feng (add refs) by using the theory of automata. In that paper, it was proved that for every Pisot number
% the gap sequence $\lambda_{n+1}-\lambda_{n}$ is the image of a substitution sequence over a finite alphabet of symbols. When considering the Pisot number $q_m$ defined as the greatest root of the polynomial 
% $$x^m-x^{m-1}-\cdots-1$$
% the substitution coding the gap sequence coincides with the \emph{Rauzy substitution over $m$ letters}  $\sigma_m$ defined by
% % For instance, the first iterations of $\sigma_m$ on $1$ are for $m=2$ (and hence $q$ equal to the golden ratio)
% % $$
% % \sigma_2(1)=12;\quad \sigma_2^2(1)=\sigma_2(12)=\sigma_2(1)\sigma_2(2)=121,\quad \sigma_2^3(1)=12112$$
% % Defining $d_m(1)=1$ and $d_m(k)=q^{k-1}-\sum_{i=0}^{k-1}q^{-i}$ we have
% % $$\lambda_{n+1}-\lambda_n=d(\mathbf v_{m,k})$$
% % where $\mathbf v_{m}:=\lim_{n\to\infty} \sigma_m^n(1).$

% \medskip

% \noindent

% In this paper we consider the spectra
% \[
% \Lambda_m := \Lambda(q_m)
% \]
%  e we provide explicit and computable lower estimates of the form 
% \begin{equation}\label{e1}
% \lambda_{k+N}-\lambda_{k}\geq N\gamma_N \end{equation}
% for all $N\in\NN$, $k\in\ZZ$, and some $\gamma_N>0$, that provide a description of the local structure of the spectrum. The analysis of such problem was addressed from a number theoretic perspective, as well as from a symbolic dynamical one -- for an overview on the symbolic dynamical aspects of the present paper see the seminal paper \cite{seminalrauzy}, the books \cite{Loth,substitutionberthebook,queffel} and, for more specific results concerning our work \cite{balancetrib, generalbalance}. It is well known that $\Lambda_m$ admits a finite upper density,  and the paper \cite{adam} proves that $\Lambda_m$ has a bounded distance to a lattice. This implies the existence of uniform linear lower bounds for the gaps. Such bounds are explicited via asymptotic estimates  in \cite{bal}. The novelty of the present analysis relies on explicit, computable formulas for all $N$ and a direct connection with Fibonacci number systems. 





% \smallskip
% The paper is organized as follows. In Section \ref{sec2} we recall the construction $m$-bonacci chains and the main properties of the  underlying namely the \emph{$m$-bonacci words}. Such bi-infinite words are built upon the celebrated Rauzy substitutions, that have a preeminent role in symbolic dynamics and combinatorics on words \cite{Loth, br10}. 
% In Section \ref{sec4} we establish explicit gap conditions \eqref{e1} in the general case,


% Our methodology is maily based on theoretical tools from discrete mathematics, combinatorics on words, and diophantine approximation, see for instance the books \cite{meyer, Loth}. In particular, the $m$-bonacci chains belong to the wide framework devoted to quasicrystals generated by unimodular Pisot substitutions, spectral and combinatorial properties of $m$-bonacci chains can be found in [cita,cita,cita]. 
% Here, we are interested in weak gap conditions on the $m$-bonacci chains, namely in estimates of the form 
% $$\lambda_{k+N}-\lambda_{k}\geq N\gamma_N$$
% for all $N,k\in\NN$, some $\gamma_N>0$ and with $\lambda_k$ belonging to some $m$-bonacci chain $\Lambda_m$. 
% The interest in the class of lower bounds is motivated by the following results in the literature. 


% use the notion of \emph{upper density}, that requires a careful analysis of the combinatorial properties of the $m$-bonacci chains. Our results  offer a detailed analysis of the $m$-bonacci chains with general order $m$, their upper densities and their gap conditions, by extending classic results in nonharmonic analysis. In particular, this study allows us to use a generalization of Ingham's inequalities  \cite{Ingham}, that is the well known Beurling  theorem \cite{beu}. 


% We consider sums of the form $$f(t):=\sum_{k=-\infty}^\infty a_k e^{i\lambda_k t}$$
% where $a_k$ are complex coefficients and $\lambda_k$ belong to a $m$-bonacci chain $\Lambda_m$ for some $m\geq 2$. 





% Our main result relies on a classic Beurling  theorem \cite{beu} and on the explicit computation of the upper density $D^+=D^+(\Lambda_m)$ of $\Lambda_m$ (see \cite{pol}) which requires a careful investigation of the combinatorial properties $\Lambda_m$. In particular, our arguments rely on tools from symbolic dynamics, discrete combinatorics and number theory. 

% Next we focus on the particular cases of the \emph{Fibonacci quasicrystal}, corresponding to the case $m=2$, and the \emph{Tribonacci quasicrystal}, corresponding to $m=3$. For these sets, we are able to establish explicit weakened gap conditions.

% that are the base of the following results.
% \begin{theorem}
% The Fibonacci quasicrystal $\Lambda_2=(\lambda_k)_{k\in\ZZ}$ satisfies the following gap condition for all $N\geq 1$ 
% $$\lambda_{N+k}-\lambda_k\geq N\gamma_{2,N}\qquad \text{for all }k\in\ZZ$$
% where 
% $$\gamma_{2,N}:=\frac{1}{\varphi^2}+\frac{1}{\varphi^4}-\frac{1}{N\varphi^3}.$$
% Moreover, for every $N$ and for every bounded interval $I$ of length 
% $|I|>2\pi/ \gamma_{2,N}$
% there exist two positive constants  $c_1,c_2>0$ such that every finite sum of the form 
% $$f(t)=\sum_{k=-\infty}^\infty a_k e^{i\lambda_k t}\qquad a_k\in\CC$$
% (with a finite number of $a_k$ different from $0$)  
% satisfies
% \begin{equation}
% c_1\sum_{k=-\infty}^{+\infty}|a_k|^2\leq \int_I |f(t)|^2dt\leq c_2 \sum_{k=-\infty}^{+\infty}|a_k|^2.
% \end{equation}  
% \end{theorem}
% The estimate on $\gamma_{2,N}$ is refined in Proposition \ref{2gaptrib}. The proof relies on a theorem on weakened gap conditions for Ingham-type inequality \cite{bkl02} and on combinatorial properties of the Fibonacci word, namely the binary word underlying the construction of the Fibonacci quasicrystal. 

% The next result concerns the Tribonacci chain. We recall that the Tribonacci sequence $(T_h)$ is recursively defined as $T_0=T_1:=0$,  $T_2:=1$ and $T_{h+3}:=T_h+T_{h+1}+T_{h+2}$.
% \begin{theorem}
% The Tribonacci chain $\Lambda_3=(\lambda_k)_{k\in\ZZ}$ satisfies the following gap condition for all $N\geq 1$ 
% $$\lambda_{N+k}-\lambda_k\geq N\gamma_{3,N}\qquad \text{for all }k\in\ZZ$$
% where 
% $$\gamma_{3,N}:=\frac{1}{N}\left(\sum_{h=3}^L x_h\left(\frac{T_{h-1}}{\tau}+\frac{T_{h-2}}{\tau^2}+\frac{T_{h-3}}{\tau^3}\right)-2\right)$$
% where $x_3\cdots x_L$ is the Tribonacci expansion of $N$, namely a binary sequence satisfying $N=\sum_{h=3}^N x_hT_h$, where $(T_h)$ is the Tribonacci sequence. 

% Moreover, for every $N$ and for every bounded interval $I$ of length 
% $|I|>2\pi/ \gamma_{3,N}$
% there exist two positive constants  $c_1,c_2>0$ such that every finite sum of the form 
% $$f(t)=\sum_{k=-\infty}^\infty a_k e^{i\lambda_k t}\qquad a_k\in\CC$$
% (with a finite number of $a_k$ different from $0$)  
% satisfies
% \begin{equation}
% c_1\sum_{k=-\infty}^{+\infty}|a_k|^2\leq \int_I |f(t)|^2dt\leq c_2 \sum_{k=-\infty}^{+\infty}|a_k|^2.
% \end{equation}  
% \end{theorem}



% Symbolic dynamics provides an extensive range of theoretical tools in this area; among them, we shall focus on a key point: the digit frequency of the fixed point $\mathbf v_m$ of the Rauzy substitution $\sigma_m$ is related to the eigenvalues and eigenvectors of a matrix (the \emph{incidence matrix}) describing the Rauzy substitution itself. In particular, if the largest eigenvalue of the incidence matrix (called the \emph{Perron-Frobenius eigenvalue}) has some algebraic properties (showed in Proposition \ref{p11}) then we can explicitly compute the digit frequencies of the words $\mathbf v_m$ (Proposition \ref{p22}) which will allow us, in Section \ref{sec3}, to get the uppper density of  $\Lambda_m$.

%  Let $\mathcal A=\{a_0,\dots,a_m\}$ be a finite alphabet, whose elements are called in the following \emph{digits}, and let $\mathcal A^*$ be the set of finite words over $\mathcal A$. A \emph{substitution} is a function $\sigma: \mathcal A^*\to\mathcal A^*$ such that $\sigma(uv)=\sigma(u)\sigma(v)$ for all words $u,v\in \mathcal A^*$ and such that $\sigma(a)$ is not the empty word for every letter $a\in\mathcal A$.

% The domain $\mathcal A^*$ of a substitution $\sigma$ can naturally be extended to the set of infinite sequences $\mathcal A^\NN$ and to set of bi-infinite sequences $\mathcal A^\ZZ$ by concatenation. In particular if $\mathbf w=w_1w_2\cdots \in \mathcal A^\NN$ then $\sigma(\mathbf w)=\sigma(w_1)\sigma(w_2)\cdots$ and if $\mathbf w=\cdots w_{-1}.\,w_0w_1\cdots \in \mathcal A^\ZZ$ then $\sigma(\mathbf w)=\cdots \sigma(w_{-1}).\,\sigma(w_0)\sigma(w_1)\cdots$. 

% \begin{definition}[Rauzy substitution and $m$-bonacci word]
%     Let $\mathcal A_m:=\{1,\dots,m\}$ with $m\geq 2$ and $\mathcal A_m^*$ be the set of finite words with digits in $\mathcal A_m$. For every $m\geq 2$, we consider \emph{Rauzy substitution} $\sigma_m:\mathcal A_m^*\to \mathcal A_m^*$ given by
% $k\mapsto 1(k+1)$ for $k=1,\dots,m-1$ and $m\mapsto 1$. 

% For all $m\geq 2$ let
%     $$\mathbf v^+_m:=\lim_{n\to \infty}\sigma^n_m(1).$$
    
%    The words $\mathbf v^+_2=\sigma_2^\infty(1)=12112121\cdots$  and $\mathbf v^+_3=\sigma_3^\infty(1)=1213121121312\cdots$ are respectively called the \emph{Fibonacci word\footnote{Actually, the Fibonacci word is formally defined as the binary word obtained as the fixed point of $\sigma_2$ on the alphabet $\{0,1\}$ rather than $\{1,2\}$. Clearly, this substitution in the alphabet does not affect the combinatorial properties of the Fibonacci word and it is commonly adopted in the more general framework of the Rauzy substitutions.}} and \emph{Tribonacci word}. For $m> 3$, the word $\mathbf v^+_m$ is called the \emph{$m$-bonacci word}. 

% The \emph{$m$-bonacci bi-infinite word} is defined by
% $$\mathbf v_m:=\lim_{n\to \infty}\sigma^{nm}_m(1.1).$$
% \end{definition}
% \begin{remark}
% Note that the definition of $\mathbf v^+_m$ is simply the limit of $\sigma_m^n(1)$ as $n\to \infty$ (see \cite{Loth} for its existence), whereas the definition of the digits with negative indexes of $\mathbf v_m$ relies on the subsequence $(mn)_{n\in\NN}$ because there is no 
% $\lim_{n\to \infty} \sigma_m^{n}(1.1)$ ---- see \cite[Lemma 4.3]{aperiodic}.\end{remark}
 
% The bi-infinite word encodes the alternance of gaps between consecutive elements of $\Lambda_m$. The size of each gap is given by following construction, see for instance \cite[Section 4]{aperiodic} for a general introduction, and \cite{baake}   for the specific case of $m$-bonacci chains. 
% Let $\rho_m$ be \emph{Perron Frobenius} (PF) eigenvalue of  the incidence matrix $M_m$ of $\sigma_m$ 
% $$M_{m} = 
% \begin{pmatrix}
% 1 & 1 & 1 & \cdots & 1 & 1 \\
% 1 & 0 & 0 & \cdots & 0 & 0 \\
% 0 & 1 & 0 & \cdots & 0 & 0 \\
% 0 & 0 & 1 & \cdots & 0 & 0 \\
% \vdots & \vdots & \vdots & \ddots & \vdots & \vdots \\
% 0 & 0 & 0 & \cdots & 1 & 0
% \end{pmatrix}.$$

%  In particular $\rho_m$ is the  greatest root of $x^m-x^{m-1}-\cdots-1$. 
% $$d_m: j\mapsto \sum_{h=0}^{m-j}\rho_m^{-h} \quad j=1,\dots,m.$$
% Note that $(d_m(1),\dots, d_m(m))$ is a left eigenvector of $M_n$ corresponding to $\rho_m$.  
% The map is normalized so that the smallest gap is equal to $1$. In particular, $d_m(1)=\rho_m$ and $d_m(m)=1$. 
% For $m\geq 0$ consider the \emph{gap sequence} $(g_{m,k})_{k\in\ZZ}$ such that $g_{m,0}:=0$ and 
% $$g_{m,k}:= d_m(\mathbf v_{m,k})\qquad \text{for } k\in\ZZ,\, k\not=0$$
% namely if the $k$-th digit of $\mathbf v_m$ is $j$ then $g_{m,k}$ is $d(j)=1+\rho_m^{-1}+\dots+\rho_m^{m-j}$. Finally we are in position to define $\Lambda_m$ 
% \begin{definition}[$m$-bonacci chains]
% For $m\geq 2$, the \emph{$m$-bonacci chain} is
% $$\Lambda_m:=\left\{\sum_{h=0}^k g_{m,h} , \,\,-\sum_{h=0}^k g_{m,-h} \mid k\in \NN\right\}.$$
% \end{definition}
% \begin{remark}[Some properties of $\mathbf v_m$ and $\Lambda_m$]
% The digit $j$ in the bi-infinite word $\mathbf v_m$ has frequency $\rho_m^{-j}$. From this, one can deduce that the density $D_m$  of $\Lambda_m$  exists and it is equal to $\beta_m$. Moreover $D_m=\beta_m\to 2$ as $m\to \infty$ \cite{aperiodic}.

% The gaps in the set $\Lambda_m$ are tailored to guarantee  $\Lambda_m$ to be a Meyer set (a class of mathematical quasicrystals), namely both $\Lambda_m$ and $\Lambda_m-\Lambda_m$ are discrete and relatively dense sets \cite{meyer}. 

% \end{remark}






% % \begin{definition}[Uniform frequency]
% % Let $\mathbf w$ be an infinite (bi-infinite) word and let $i=1,\dots,m$. If for all $k\geq 0$ ($k\in \ZZ$) the limit
% % $$\lim_{n\to\infty}\frac{|w_{k+1}\cdots w_{k+n}|_j}{n}$$
% % exists uniformly with respect to $k$, then it is called the \emph{uniform frequency $f_j(\mathbf w)$ of the digit $a_j$ in $\mathbf w$}. Equivalently, if above limit exists, we can define
% % $$f_j(\mathbf w):=\lim_{n\to \infty} \frac{\sup\{|w|_j\mid w \text{ is a subword of length $n$ of $\mathbf w$}\}}{n}.
% % $$
% % \end{definition}





% % CORRETTO FINO A QUI

% % We point out that the incidence matrices of the Fibonacci and Tribonacci substitutions are respectively 
% % $$M_{\sigma_2}=\begin{pmatrix}
% % 1&1\\
% % 1&0
% %    \end{pmatrix}; \qquad M_{\sigma_3}=\begin{pmatrix}
% % 1&1&0\\
% % 1&0&1\\
% % 1&0&0
% %    \end{pmatrix};$$ 
% %   whereas for a generic Rauzy substitution $\sigma_m$  with $m\geq 4$ we have
% % $$M_{\sigma_m}=\begin{pmatrix}
% %     1&1&0&\cdots&0\\
% %     1&0&1&\cdots&0\\
% %     \vdots&\vdots&\ddots&\ddots&\vdots\\
% %     1&0&\cdots&0&1\\
% %     1&0&\cdots&0&0
% % \end{pmatrix}$$
% % Note the recursive relation
% % \begin{equation}\label{Mrec}M_{\sigma_m}=\begin{pmatrix}
% %   M_{\sigma_{m-1}}&\mathbf e^{T}_{m-1,m-1}\\
% %     \mathbf e_{m-1,1}&0\\
% % \end{pmatrix}\end{equation}
% % where $\mathbf e_{m,j}$ is the $j$-th element of the canonical base of $\RR^m$. 



% % \begin{definition}
% % The \emph{Perron-Frobenious eigenvalue} of a substitution is defined as the largest eigenvalue of its adjacency matrix. We denote by $\rho_m$ the Perron-Frobenious eigenvalue of $\sigma_m$. 

% % A substitution $\sigma$ is a \emph{Pisot substitution} if its Perron-Frobenious eigenvalue is a Pisot number, namely an algebraic integer greater than $1$ whose conjugates are less than $1$ in modulus. 
% % \end{definition}
% % By a direct computation we have that, for the case of Fibonacci and Tribonacci substitution we respectively have that $\rho_2$  is the Golden Mean, namely the greatest solution of $\lambda^2=\lambda+1$, and $\rho_3=1.83929...$ is the Tribonacci constant, namely the greatest solution of $\lambda^3=\lambda^2+\lambda+1$. Both $\rho_2$ and $\rho_3$ are {Pisot numbers}, and consequently both $\sigma_2$ and $\sigma_3$ are Pisot substitutions. 
% % This is a general fact that we recall below. 


% % \begin{proposition}\label{p11}
% %     For all $m\geq 2$, $\sigma_m$ is a Pisot substitution. 
    
% %     In particular, the {Perron-Frobenious} eigevanlue $\rho_m$ of $\sigma_m$ is the Pisot number whose minimal polynomial  $\lambda^m-\lambda^{m-1}-\cdots-\lambda-1$. Moreover the vector
% %     $$d_m:=(\rho_m^{-1},\dots,\rho_m^{-m}) $$
% %     is a left eigenvector associated to $\rho_m$ whose $\ell_0$ norm is equal to $1$. 
% % \end{proposition}
% %      \begin{proof}
% %      We proceed by first proving that for all $m\geq 2$ the characteristic polynomial of $M_{\sigma_m}$ is $P_m(\lambda)=(-1)^m(\lambda^m-\lambda^{m-1}-\dots-1)$. The case $m=2$ is obtained by a direct computation and we assume it as an inductive hypothesis. Now, remark that in general, setting $I_m$ the $m\times m$ identity matrix, it follows from \eqref{Mrec}
% % that      \begin{align*}M_{\sigma_m}-\lambda I_m&=\begin{pmatrix}
% %         M_{\sigma_{m-1}}-\lambda I_{m-1}&\mathbf e_{m-1,m-1}^T\\
% %        \mathbf e_{m-1,1}&-\lambda\\
% %      \end{pmatrix}     =\begin{pmatrix}
% %        (1^m)^T& \mathbf T_{m-1}\\
% %        1&-\lambda \mathbf e_{m-1,m-1}\end{pmatrix}\end{align*}
% %     where $\mathbf T_{m-1}$ is a lower triangular matrix whose diagonal elements are $1$, consequently its determinant is $1$. Therefore, also by using the inductive hypothesis, 
% %      \begin{align*}
% %          P_m(\lambda)&=\det (M_{\sigma_m}-\lambda I_m)\\
% %          &=(-1)^{m+1}\det(\mathbf T_{m-1})+(-1)^{2m}(-\lambda)\det (M_{\sigma_{m-1}}-\lambda I_{m-1}))\\
% %          &=(-1)^{m}(-1+(-1)^{m+1}\lambda P_{m-1}(\lambda))\\
% %          &=(-1)^{m}(-1+(-1)^{2m}\lambda (\lambda^{m-1}-\lambda_{m-2}-\cdots-1))\\
% %          &=(-1)^m(\lambda^m-\lambda^{m-1}-\cdots-1). 
% %      \end{align*}
% %      The fact that $\rho_m$ is a Pisot hence follows by  \cite{br} -- see also \cite{fr}.  It is left to check that $d_m$ is a left eigenvector. Now, we have that the vector equality 
% %      $$d_m M_{\sigma_m}=\rho_md_m^T$$
% %      is equivalent to $$\left(\sum_{j=1}^m \rho_{m}^{-j},\rho_m^{-1},\dots,\rho_m^{-m+1}\right)=(1,\rho^{-1},\dots,\rho_m^{m-1})$$
% % which follows from the fact that $\rho^{m}_m=\sum_{j=0}^{m-1}\rho^{j}_m.$
% % \end{proof}
% % Proposition \ref{p1} ensures to the Rauzy substitutions several properties, that are typical of Pisot substitutions and that we recall below from the book \cite{br10} -- see in particular Chapter 1 and Chapter 5.
% % \begin{proposition}[Rauzy words and their frequencies]\label{p22}
% %     For all $m\geq 2$ let 
% %     $$\mathbf w_m:=\lim_{n\to \infty}\sigma^n_m(1).$$
% %     Then $\mathbf w_m$ is a uniformly recurrent infinite word, namely every factor of  $\mathbf w_m$ occurs in $\mathbf w_m$ infinitely often and the gap between one occurrence and the other is uniformly bounded.

% %     Moreover for all $j$, $\mathbf w_m$ admits uniform frequency $f_{\mathbf w_m}(j)$ and it is equal to the $j$-th component of the left eigenvector $d_m$, that is 
% %     $f_{\mathbf w_m}(j)=\rho_m^{-j}$.
% % \end{proposition}
 
% We conclude this section with a standard result relating the $m$-bonacci chains and $m$-bonacci numbers, whose proof is reported for clarity purposes. 



% \begin{proposition}\label{p0}
% Setting $F_{-1}^{(m)}=1$ and $F_{-k}^{(m)}=0$ for all $k> 1$, we have all $h\geq 0$ and for all $m\geq 2$ 
% $$\rho_m^h=(d_m(1),\dots, d_m(m))\cdot (F_{h-1}^{(m)},\dots, F_{h-m}^{(m)})$$
% \end{proposition}
% \begin{proof}
% By construction for all $h\geq 0$ we have 
% $$(F_{h}^{(m)},\dots, F_{h+1-m}^{(m)})^T=M_n (F_{h-1}^{(m)},\dots, F_{h-m}^{(m)})^T.$$
% As remarked above $(d(1),\dots,d(m))$ the left PF eigenvector of $M_n$. Then 
% \begin{align*}P_{h,m}&:=(d_m(1),\dots, d_m (m))\cdot(F_{h}^{(m)},\dots, F_{h+1-m}^{(m)})\\
% &=\rho_m P_{h-1,m}=\rho^{h+1}_m P_{-1,m}=\rho_m^{h+1}(1,0,\dots,0)\\
% &=\rho^{h+1}_m.\end{align*}
% \end{proof}
% We extend the $m$-bonacci word $\mathbf w_m$ to the
% The ‘positive' part of $\mathbf v_m$  coincides with $\mathbf w_m$, and $\sigma^{nm}_m(1.1)=\sigma_m^{nm}(1).\sigma_m^{nm}(1)$ for all $n$. Therefore the results of Proposition \ref{p22} also apply to $\mathbf v_m$ which is a uniformly recurrent bi-infinite sequence whose uniform frequencies satisfy $f_{\mathbf v_m}(j)=\rho_m^{-j}$ for all $j=1,\dots,m$. 





% \subsection{The $m$-bonacci chains}\label{sfib}\label{sec3}
% In the previous section we introduced the $m$-bonacci bi-infinite word $\mathbf v_m=\cdots v_{-1}.v_{0}v_1\cdots$, that is a bi-infinite extension of the fixed point $\mathbf w_m$ of the Rauzy substitution $\sigma_m$. The digits of $\mathbf v_m$ belong to the alphabet $\mathcal A_m=\{1,\ldots,m\}$, and for all $j=1,\ldots,m$ there exists the uniform frequency of the digit $j$. It satisfies
% \begin{equation}\label{frecap}
% f_{\mathbf v_m}(j)=\rho_m^{-j}\end{equation}
% where $\rho_m$ is the greatest real root of the polynomial $P_m(x)=x^{m}-x^{m-1}-\cdots-x-1$. We use $\mathbf v_m$ to construct the \emph{gap sequence}
% $$d_{m,k}:=f_{\mathbf v_m}(v_k)=\rho_m^{-v_k}\qquad \text{for all $k\in\ZZ$}$$
% and the sequence 
% $$\lambda_{m,0}:=0;\qquad \lambda_ {m,n}:=\sum_{k=1}^n d_{m,k},$$ $$ \lambda_{m,-n}:=-\sum_{k=1}^n d_{m,-k}\,\quad n\in\NN.$$
% \begin{definition}[$m$-bonacci chains]
% For $m\geq 2$, the \emph{$m$-bonacci chain} is
% $$\Lambda_m:=\{\lambda_{m,k}\mid k\in \ZZ\}.$$
% \end{definition}
% The set $\Lambda_m$ is a relatively dense, discrete set (namely it is a \emph{Delone set}) and it has a pure point diffraction pattern, that is, $\Lambda_m$ is a quasicrystal \cite{rauzyqc}. In the case $m=2$, $\Lambda_2$ is the celebrated \emph{Fibonacci quasicrystal}. 



% \begin{remark}[Upper density]
% The \emph{upper density} $D^+(\Lambda)$ of a discrete set $\Lambda\subset \RR$  is defined as 
% $$D^+(\Lambda):=\limsup_{r\to+\infty} \frac{\mathbf n(r)}{r}$$ where $\mathbf n(r)$ is the greatest number of elements of $\Lambda$ in an interval of length $r$. \end{remark}
% Our main result is the following. 
% \begin{theorem}
% \label{thm1}
% Let $m\geq 2$ and let $\Lambda_m=(\lambda_k)_{k\in \ZZ}$ be the corresponding $m$-bonacci chain. Then
% \begin{equation}\label{upden}D^+(\Lambda_m)=\frac{\rho^{2m}_m}{1+\rho_m^2+\dots+\rho_m^{2(m-1)}},\end{equation}
% where $\rho_m$ is the greatest root of $x^m-x^{m-1}-\cdots-x-1$.
% Moreover $D^+(\Lambda_m)\nearrow 3$ as $m\to +\infty$ and 
% \begin{equation}\label{Dapprox}
%     D^+(\Lambda_m)\geq 3-\frac{8}{2^m}\left(1-\frac{1}{2^{m+1}}\right) \quad \text{ for all } m\geq 2.
% \end{equation}

% For every bounded interval $I$ of length 
% $$|I|>D^+(\Lambda_m)=2\pi \frac{\rho^{2m}_m}{1+\rho_m^2+\dots+\rho_m^{2(m-1)}}$$
% there exist two positive constants  $c_1,c_2>0$ such that every sum of the form 
% $$f(t)=\sum_{k=-\infty}^\infty a_k e^{i\lambda_k t}\qquad a_k\in\CC$$
% satisfies
% \begin{equation}
% c_1\sum_{k=-\infty}^{+\infty}|a_k|^2\leq \int_I |f(t)|^2dt\leq c_2 \sum_{k=-\infty}^{+\infty}|a_k|^2.
% \end{equation}
% The above inequalities can fail when $|I|<2\pi D^+(\Lambda_m)$. 
% \end{theorem}

% \begin{proof}
% We fix $m\geq 2$ and we set $\lambda_k=\lambda_{m,k}$  and $d_k=d_{m,k}$, in order to enlighten the notation. 
% For every $n\geq 1$, respectively define 
% $$r^-_n:=\min_{k\in\ZZ}\{\lambda_{k+n}-\lambda_{k}\}
% \quad \text{and}\quad
% r^+_n:=\max_{k\in\ZZ}\{\lambda_{k+n}-\lambda_{k}\}$$
% Note that both $r^+_n$ and $r^-_n$ are well defined because $\Lambda_m$ is a Delone set: in particular, since $\Lambda_m$ is discrete then $r^-_n$ is bounded, and since $\Lambda_m$ is relatively dense (that is the distance between two consecutive elements is uniformly bounded from above) then $r^+_n$ is bounded as well. 
% By definition of weight, for every $k\in\ZZ$ and for every $j=1,\dots,m$, in  the subword $v_{k+1}\cdots v_{k+n}$  of $\mathbf v_m$  there are exactly $|v_{k+1}\cdots v_{k+n}|_j$ occurrences of the digit $j$. Therefore 
% \begin{align*}\frac{\lambda_{k+n}-\lambda_{k}}{n}&=\frac{1}{n}\sum_{h=1}^n d_{k+h}=\frac{1}{n}\sum_{j=1}^n \rho_m^{-v_{k+j}}\\
% &=\frac{1}{n}\sum_{j=1}^m |v_{k+1}\cdots v_{k+n}|_j\rho_m^{-j}\to \sum_{j=1}^m f_{\mathbf v_m}(j)\rho_m^{-j} 
%  \end{align*}
% as $n$ tends to $+\infty$ uniformly with respect to $k$. This in particular implies, also in view of \eqref{frecap},  that
% $$\lim_{n\to+\infty}\frac{r^+_n}{n}=\lim_{n\to+\infty}\frac{r^-_n}{n}= \sum_{j=1}^m f_{\mathbf v_m}(j)\rho_m^{-j} =\sum_{j=1}^m \rho_m^{-2j}=: s_m$$
% as the minimal and maximal distance  between $n+1$ consecutive elements of $\Lambda_m$. 
% $\limsup_{n\to +\infty} \frac{r^-(n)}{n}=\lim_{n+\infty} $
% Finally note that for all $n$, 
% $$\mathbf n(r^-_n)\leq n+1 \leq \mathbf n(r^+_n)$$
% and, consequently, 

% \begin{align*}\frac{1}{s_m}=\lim_{n\to\infty} \frac{n+1}{r^+_n}\leq \limsup_{n\to\infty} \frac{\mathbf n(r^+_n)}{r^+_n} &= \limsup_{r\to \infty}\frac{\mathbf n(r)}{r}=\limsup_{n\to\infty} \frac{\mathbf n(r^-_n)}{r^-_n}\\
% &\leq \lim_{n\to\infty} \frac{n+1}{r^-_n}=\frac{1}{s_m}.\end{align*}
% The equality \eqref{upden} hence following by remarking that ${s_m}=\frac{1+\rho_m^2+\dots+\rho_m^{2(m-1)}}{\rho^{2m}_m}.$

% Now, in \cite{wolfram} it was proved that $\rho_m\in  (2(1-2^{-m}),2)$ for all $m$ and, using the Descartes rule of signs one can prove that $\rho_m$ is an increasing sequence. Since 
% $$D^+(\Lambda_m )=\frac{\rho^{2m}_m(\rho^2_m-1)}{\rho_m^{2m}-1}$$
% we readily get that $D^+(\Lambda_m)\nearrow 3$ as $m\to \infty$.
% Moreover we have that $D^+(\Lambda_m )\geq \rho^2_m-1.$
% By replacing $\rho_m$ with $2(1-2^{-m})$ in the right hand-side of the this inequality we obtain the estimate \eqref{Dapprox}.


% The second part of the statement is a direct application of Beurling theorem on trigonometric series \cite{beu}.

% \end{proof}
% \begin{example}[Density of the Fibonacci chain]
% For the Fibonacci chain (or Fibonacci quasicrystal) we have
% $$D^+(\Lambda_2)=
% \frac{\varphi^4}{1+\varphi^2}\simeq 1.89443,$$
% where $\varphi=\rho_2$ is the Golden Ratio. Note that the estimate \eqref{Dapprox} provides $D^+(\Lambda_2)\geq 5/4=1.25$ Using the result in \cite{wolfram} this can be refined to $D^+(\Lambda_2)\geq 81/52=1.55769.5$
% In Figure \ref{fig1} we show the numerical approximation of $D^+(\Lambda_2)$ by computing the ratio $\mathbf n(r)/r$ for large $r$. The computation of $\mathbf n(r)$ is restricted to $\{\lambda_k\}_{k=0}^N$ with large $N$, that corresponds to large prefixes of the word $\mathbf w_2$. 
% This can be justified by that fact that $\mathbf w_2$ is uniformly recurrent, so every subword in the prefix occurs uniformly in the 
% rest of the word, moreover $\mathbf w_2$ has a 'self-similar structure' preventing large variations on the tail of $\mathbf w_2$. 
% % \begin{figure}[h!]
% %     \centering
% %     \includegraphics[scale=0.5]{grafico_beurling_fib.pdf}
% %     \caption{An approximation of $D^+(\Lambda_2)$ via $\mathbf n(r)/r$ with $r\in[250,500]$.}
% %     \label{fig1}
% % \end{figure}
% \end{example}
\section{Gap conditions for the spectrum of $m$-bonacci numbers }\label{sec3}
In this section we prove our main results. We begin by proving Theorem \ref{thm1i} that is 

\begin{theorema}
Let $m\geq 2$ and $N\in \mathbb{N}$. Then there exists a constant
$\gamma_{m,N}>0$ such that
\[
\lambda_{n+N}(q_m)-\lambda_n(q_m)\geq N\,\gamma_{m,N}
\quad \text{for all } n\geq 1.
\]

The following choice of $\gamma_{m,N}$ satisfies the above inequality.
Let $(x_h)\in\{0,1\}^\infty$ be the canonical $m$-bonacci expansion of $N$ and let $b_m>0$ such that the $m$-bonacci word $\mathbf v_m$ is $b_m$-balanced.
Then
\begin{equation}\label{eqgamma}
\gamma_{m,N}
=\frac{1}{N}
\sum_{j=1}^m
\left(
\sum_{h\geq 0} x_h F^{(m)}_{h-j}
- b_m
\right)
d_m(j),
\end{equation}
where
 $d_m(j)=q_m^{j-1}-\sum_{k=0}^{j-2} q_m^k$ with $j=1,\dots,m$.
\end{theorema}
\begin{proof}
First of all we remark that, by Theorem \ref{thm2},
for all $m\geq 2$, $k,N \in \NN$
\begin{align*}\label{gapw}
    \lambda_{k+N}(q_m)-\lambda_{k}(q_m)&=
    \sum_{h=1}^N(\lambda_{k+h}(q_m)-\lambda_{k+h-1}(q_m))\\
    &=\sum_{j=1}^m |\{h\mid \lambda_{k+h}(q_m)-\lambda_{k+h-1}=d_m(j),\,h=1,\dots,N\}|d_m(j)\\
    &=\sum_{j=1}^m |\mathbf v_{k+1}\cdots \mathbf v_{k+N}|_j d_m(j).
\end{align*}
As $\mathbf v_m$ is a balanced word it admits a lowest balance constant $b_m$, see Theorem \ref{thm3}. We then deduce from the above equality the lower estimate 
\begin{equation}\label{gapw}
    \lambda_{k+N}(q_m)-\lambda_{k}(q_m) \geq \sum_{j=1}^m\left(|\mathbf v_{m,1}\cdots \mathbf v_{m,N}|_j-b_m\right)d_m(j)\end{equation}
for all $m\geq 2$, $k,N \in \NN$. 


Let $(x_h)$ be the canonical $m$-bonacci expansion of $N$.
By the converse statement of Lemma \ref{l1bal}, there exists a
prefix $v$ of $\mathbf v_m$ such that

\[
(|v|_1,\dots,|v|_m)^T
=
\sum_{h\ge0}x_h M_m^h(1,0,\dots,0)^T .
\]

Now, a simple inductive argument (already remarked in \cite{generalbalance} with a slightly different notation) shows that if $M_m$ is the incidence matrix of $\sigma_m$, defined in \eqref{Mm} then
$$M^h(1,0,\dots,0)^T=(F^{(m)}_{h-1},F^{(m)}_{h-2},\dots, F^{(m)}_{h-m})^T$$
where $F^{(m)}_{-1}=1$ and $F^{(m)}_{-k}=0$ for all $k> 1$.
Then we obtain
\[
|v|_j
=
\sum_{h\ge0}x_hF_{h-j}^{(m)}
\qquad (j=1,\dots,m).
\]

Moreover,

\[
|v|
=
\sum_{j=1}^m|v|_j
=
\sum_{h\ge0}x_h
\sum_{j=1}^mF_{h-j}^{(m)}
=
\sum_{h\ge0}x_hF_h^{(m)}
=
N.
\]

Therefore $v$ is the prefix
$\mathbf v_{m,1}\cdots\mathbf v_{m,N}$ and hence

\[
|\mathbf v_{m,1}\cdots\mathbf v_{m,N}|_j
=
\sum_{h\ge0}x_hF_{h-j}^{(m)}.
\]
This, together with \eqref{gapw}, concludes the proof of the theorem. 



% This, together with Lemma \ref{l1bal}, implies that that there exists some binary finite sequence $(y_h)$ satisfying
% \begin{equation}\label{nj}N_j:=|\mathbf v_{m,1}\cdots \mathbf v_{m,N}|_j=\sum_{h=0}^\infty y_h F_{h-j}^{(m)} \qquad \text{for all } j=1,\dots,m.\end{equation}
% Then we note that 
% $$N=\sum_{j=1}^m|\mathbf v_{m,1}\cdots \mathbf v_{m,N}|_j=\sum_{h=0}^\infty \sum_{j=1}^m y_h F_{h-j}^{(m)}=\sum_{h=0}^\infty y_h F_{h}^{(m)} \qquad \text{for all } j=1,\dots,m.$$
% From this we deduce that the sequence $(y_h)$ posited in Lemma \ref{l1bal} is actually a $m$-bonacci expansion of $N$. The second part in the statement of Lemma \ref{l1} states that every choice of $(y_h)$ yields a prefix of $\mathbf v_m$ satisfying \eqref{prefix} and consequently, \eqref{nj}, then we can choose $y_h$ as the canonical expansion of $N$.


\end{proof}

% consider $m\geq 2$ and we  order $\Lambda_m$ via the strictly increasing sequence $(\lambda_{m,k})$. Our goal is to establish inequalities of the form 
% $$\lambda_{m,k+N}-\lambda_{m,k}\geq N \gamma_{m,N}\quad  \text{for all }k\in\ZZ$$
% for some positive constant $\gamma_{m,N}$, for  all $N\in \NN$.

% and for $m=2,3$, namely in the Fibonacci and the Tribonacci case, respectively. 
% First of all, we relate the gaps in $\Lambda(q_m)$ to the digit distribution of the $m$-bonacci word $\mathbf v_m$ via the following identity relation
% \begin{lemma} \label{p1} 
% For all $m\geq 2$, $k,N \in \NN$

% .
% \end{lemma}
% \begin{proof} The identity \eqref{gapw} follows from the definition of $\Lambda_m$ and states that if we want to measure the length of $[\lambda_{k},\lambda_{k+N}]$ we can count how many tiles of the form $[0,d(j)]$ for each $j=1,\dots,m$ are used to pave the $[\lambda_{k},\lambda_{k+N}]$ (recall indeed that $\lambda_{k+1}-\lambda_k=d(j)$ for some $j=1,\dots,m$). The number of such tiles is provided exactly by the number of occurrences of $j$ in $v_{k+1}\cdots v_{k+N}$, namely $|v_{k+1}\cdots v_{k+N}|_j$. 
% \end{proof}
% % Below, we use this argument by distinguishing the cases $m=2,3$ and general $m>3$. 
% The next result allows to restrict ourselves to investigation of the digit distribution to  the first terms of $\mathbf v^+_m$. 
% \begin{lemma} \label{p2}
% For all $m\geq 2$ there exists positive a constant $b_m$ such that for every $j=1,\dots,m$, $k\in\ZZ$ and $N\geq 1$
% $$|\mathbf v_{m,k+1}\cdots \mathbf v_{m,k+N}|_j\geq  |\mathbf v_{m,1}\cdots \mathbf v_{m,N}|_j-b_m.$$
% In particular if $b_2=1, b_3=2$ and in general $b_m\leq m-1$ for all sufficiently large $m$.
% \end{lemma}
% \begin{proof}
% The claim of follows by the more general fact that $\mathbf v_m$ is a $m$-balanced word, namely if $u,w$ are subwords of the same length of $\mathbf v_m$ then for all $j=1,\dots,m$ we have $||u|_j-|w|_j|\leq b_m$ , for some constant,  see \cite{Loth} for a general introduction and the case $m=2$, \cite{balancetrib} for the case $m=3$ and \cite{generalbalance} for the general case. 
% \end{proof}
% Recall from \cite{mnacci} that every $N\in\NN$ admits an expansion of the form 
% $$N=\sum_{h=0}^{\infty}x_h F^{(m)}_h\qquad x_h\in\{0,1\}$$
% with all $x_h=0$ except for finitely many. 

% of the present 

% so we we refer to $(x_h)$ as  \emph{a} $m$-bonacci expansion.
% we mean  the  lexicographic greatest (or greedy) expansion of $N$. 

% The next result relates the first digits of $\mathbf v_m$ to the $m$-nacci expansions. 

% \begin{theorem}
% %[Gap conditions ]
% \label{thm1}
% Let $m\geq 2$, $N\in \NN$. For any choice $(x_h)$ of a $m$-bonacci expansion of $N$ the constant 
% $$\gamma_{m,N}:=\frac{1}{N}\left(\sum_{j=1}^m\left(\ \sum_{h=0}^\infty (x_h F^{(m)}_{h-j})-b_m\right)d_m(j)\right)$$
% satisfies the gap condition
% \begin{equation}\label{gengaptr}
% \lambda_{m,N+k}-\lambda_{m,k}\geq N\gamma_{m,N} \quad \text{for all } k\in \ZZ,
% \end{equation}
% \end{theorem}
% \begin{proof}
% In view of Lemma \ref{p1} and Lemma \ref{p2} we have that any constant $\gamma_{m,N}$ sastisfying \eqref{gengaptr} can be chosen as  
% \begin{equation}\label{eqgamma} \gamma_{m,N}=\frac{1}{N}\sum_{j=1}^m\left(|\mathbf v_{m,1}\cdots \mathbf v_{m,N}|_j-b_m\right)d_m(j).\end{equation}
% So the rest of the proof is devoted to the computation of the digit weights $|\mathbf v_{m,1}\cdots \mathbf v_{m,N}|_j$. 



% Every $m$-bonacci expansion is finite,  in particular
% $$N=\sum_{h=1}^{H} x_h F_h^{(m)}, \text{for some } H\leq F_N^{(m)} \text{ depending on $x_h$}.$$ By \cite[Lemma 2.3] {generalbalance} (see also \cite{seminalrauzy}) and setting $\sigma_m^0(0)=\sigma_m(0)$ equal to the empty word, we have
% \begin{equation}\label{prefix} \mathbf v_{m,1}\cdots \mathbf v_{m,N}=\sigma_m^{H(N)}(x_{H(N)}) \sigma_m^{H(N)-1}(x_{H(N)-1})\cdots \sigma_m^{0}(x_0) .\end{equation} 
% This formula relates the expansions $(x_h)$ of $N$ to a factorization of $\mathbf v_{m,1}\dots \mathbf v_{m,N}$. Indeed  $\mathbf v_{m,1}\cdots \mathbf v_{m,N}$ is written as a concatenation of the images of  $(x_h)$. In particular, if $x_h=0$ then the corresponding factor $\sigma_m^h(0)$ is the empty word. If otherwise $x_h=1$ then the corresponding factor is $\sigma_m^h(1)$ whose length is $F_h^{(m)}$, as we see right below. In general,
% $$\sigma_m^h(1) = \sigma_m^{h-1}(1) \sigma_m^{h-2}(1) \cdots \sigma_m^{0}(1) $$
% $\text{for all }h\geq 1$. 
% From this, we readily have by and inductive argument, that
% $$|\sigma^h_m(1)|=F_{h}^{(m)}$$
% and that for all $j=1,\dots,m$ $$|\sigma^h_m(1)|_j=F_{h-j}^{(m)}$$
% for all $h$ -- with the convention that $F^{(m)}_{-k}=0$ for all $k\in\NN$. Recalling that $|\sigma_m^h(x)|=0=x|\sigma_m^h(1)|$ if $x=0$,   for all $j=1,\dots,m$
% \begin{align*}
%     |\sigma_m^{H }(x_H)\sigma_m^{(H-1) }(x_{H-1})\cdots \sigma_m^{0}(x_0)|_j&=\sum_{h=0}^H  |\sigma_m^{h }(x_h)|_j\\
%     &=\sum_{h=0}^H x_h F_{h-j}^{(m)}.
% \end{align*}
% Plugging the above identity in \eqref{prefix} and then in \eqref{eqgamma} concludes the proof of \eqref{gengaptr}. 

% \end{proof}


The next results discuss the cases $m=2,3$. 


\begin{corollary}[Fibonacci case]\label{cor1}
For all $N\in \NN$ we have the gap condition
\begin{equation}\label{gengapfib}
\lambda_{2,N+k}-\lambda_{2,k}\geq N\gamma_{2,N} \quad \text{for all } k\in \ZZ,
\end{equation}
where
$$\gamma_{2,N}:=\frac{1}{\varphi}+\frac{1}{\varphi^3} -\frac{1}{N}\left(\frac{1}{\varphi^4}+\varphi\right).$$
The estimate is sharpened in the case $N=2$ by setting  $$\gamma_{2,2}=\frac{\varphi}{2}.$$

\end{corollary}
\begin{proof}
Set for brevity $F^{(2)}_n=F_n$ (so that $F_0=1, F_1=2,F_2=3\cdots$, in particular $F_n$ is the $n+2$th Fibonacci number) and let $(x_h)$ be the canonical Fibonacci expansion of $N$. Note that $d_2(1)=1$ and $d_2(2)=1/\varphi$. Since the balance constant $b_2$ is equal to $1$, by Theorem \ref{thm1i} we have that 
\eqref{gengapfib} holds with 
\begin{align*}\gamma_{2,N}
&=
\frac{1}{N}
\sum_{j=1}^2
\left(
\sum_{h=0}^\infty x_h F_{h-j} - 1
\right)d_2(j)\\
&=
\frac{1}{N}
\left[
\left(\sum_{h=0}^\infty x_h F_{h-1} -1\right)
+
\left(\sum_{h=0}^\infty x_h F_{h-2} -1\right)\frac{1}{\varphi}
\right] \\
&=
\frac{1}{N}
\left(
\sum_{h=0}^\infty x_h\left(F_{h-1}+F_{h-2}\frac{1}{\varphi}\right)
-\varphi
\right).
\end{align*}
Since 
$$F_{h-1}=\frac{1}{\varphi} F_h +  \left(-\frac{1}{\varphi}\right)^{h+2};\quad F_{h-2}=\frac{1}{\varphi^2} F_h -  \left(-\frac{1}{\varphi}\right)^{h+2} $$
% for all $h$ and $\sum_{h=0}^\infty x_h F_h=N$, 
we obtain
\begin{align*}\gamma_{2,N}&=\frac{1}{ N}\left(\left(\frac{1}{\varphi}+\frac{1}{\varphi^3}\right) \sum_{h=0}^\infty x_h F_h+ \frac{1}{\varphi^4}\sum_{h=0}^\infty  \frac{x_h}{(-\varphi)^h}-\varphi\right)\\
&=\frac{1}{\varphi}+\frac{1}{\varphi^3} +\frac{1}{ N}\left(\frac{1}{\varphi^4}\sum_{h=0}^\infty \frac{x_h}{(-\varphi)^h}-\varphi\right)\\
&\geq\frac{1}{\varphi}+\frac{1}{\varphi^3}  -\frac{1}{ N}\left(\frac{1}{\varphi^3}\sum_{h=1}^\infty \frac{1}{\varphi^{2h}}+\varphi\right)\\
&= \frac{1}{\varphi}+\frac{1}{\varphi^3} -\frac{1}{N}\left(\frac{1}{\varphi^4}+\varphi\right).
\end{align*}


In the particular case $N=2$, we note that the subword $22$ does not occur
in the Fibonacci word.
Therefore,
\[
\lambda_{2,k+2}-\lambda_{2,k}\ge d_2(1)+d_2(2)=\varphi,
\]
and one can choose $\gamma_{2,2}=\varphi/2$.


\end{proof}
\begin{corollary}[ Tribonacci case]\label{cor2}
Let $N\in \NN$ and let $(x_h)$ be the canonical Tribonacci expansion of $N$. Then we have the gap condition
\begin{equation}\label{gengaptr}
\lambda_{3,N+k}-\lambda_{3,k}\geq N\gamma_{3,N} \quad \text{for all } k\in \ZZ,
\end{equation}
where
\begin{align*}\gamma_{3,N}&:=\frac{1}{N}\left(\sum_{h=0}^\infty  x_h\tau^{h}-2(\tau^2-1)\right)\\&\geq \frac{1}{C_1\tau^3} -\frac{1}{N}\left({\frac{1}{2C_1\tau^3}\log_\tau \left( \frac{N +\frac{1}{2}}{C_1} \right)+2\tau^2-2-\frac{1}{C_1\tau^3})}\right).\end{align*}
where $\tau$ is the Tribonacci constant (i.e. the greatest root of $\tau^3-\tau^2-\tau-1$) and  $$C_1:= \frac{\tau^3}{4\tau^2 - 3\tau - 3}\sim 0.18. $$

The estimate is sharpened in the case $N=2$ by setting  $$\gamma_{3,2}=\frac{\tau^2-\tau}{2}.$$
\end{corollary}

\begin{proof}
Set for brevity $F^{(3)}_n=T_n$ (so that $T_0=1, T_1=2,T_2=4,T_3=7 \cdots$, in particular $T_n$ is the $n+3$th Tribonacci number) and let $(x_h)$ be the canonical $m$-bonacci (i.e., Tribonacci) expansion of $N$. Note that $b_3=2$ and

$$d_2(1)=1\,\quad d_2(2)= \tau-1=\frac{1}{\tau}+\frac{1}{\tau^2},   \quad d_2(3)=\tau^2-\tau-1=\frac{1}{\tau}.$$

By Theorem \ref{thm1i} we have that \eqref{gengaptr} holds with 
\begin{equation}\label{r}\begin{split}\gamma_{3,N}=&\frac{1}{N}\sum_{h=0}^\infty x_h\left(T_{h-1}+\frac{\tau+1}{\tau^2}T_{h-2}+\frac{1}{\tau}T_{h-3}\right)-2\frac{\tau^2-1}{N}.\end{split}\end{equation}
One can check by induction  that 
\begin{equation}\label{rec} \tau^h= T_{h-1}+\frac{\tau+1}{\tau^2} T_{h-2}+ \frac{1}{\tau}T_{h-3}.\end{equation}
Replacing the above expression in \eqref{r}  we have that 
\begin{align}\gamma_{3,N}=\frac{1}{N}\left( \sum_{h=0}^\infty x_h \tau^h-2(\tau^2-1)\right).\label{g}
\end{align}
By \cite{binet} and by the fact  that $T_h$ is the $h+3$-th Tribonacci number, we have  \begin{equation}\label{binet} |T_{h}-C_1\tau^{h+3}|\leq \frac{1}{2} \end{equation}
for all $h$ and, consequently, the lower estimate
$$\tau^h\geq \frac{1}{C_1\tau^3}\left(T_h-\frac{1}{2}\right)$$
yielding
\begin{align}\gamma_{3,N}\geq\frac{1}{C_1\tau^3}-\frac{1}{N}\left( \frac{1}{2C_1\tau^3}\sum_{h=0}^\infty x_h +2(\tau^2-1)\right).\label{g1}
\end{align}
To estimate the sum $\sum_{h=0}^\infty x_h$, let $H=H(N)$ be such that $N\in[T_H,T_{H+1})$ so that $x_h=0$ for all $h>H$.  
Then $N\geq T_{H}$ and by \eqref{binet} we deduce
$$N\geq C_1\tau^{H+3} - \frac{1}{2}$$
and
$$\sum_{h=0}^\infty x_h\leq H+1 \le \log_\tau\left(\frac{N+1/2}{C_1}\right)-2$$
Plugging the above inequality in \eqref{g1} we conclude
% By standard results on the sums of the Tribonacci numbers and by \eqref{binet} we 
% \begin{align*}N&\leq \sum_{h=0}^{H(N)}T_h=
% \frac{T_{H+2} + T_H - 3}{2}\\&\leq 
% \frac{C_1\tau^{H+3}(\tau^{2}+1) }{2}-1\end{align*}
% and consequently
% $$H(N) \le \log_\tau\left(\frac{N+0.5}{C_1}\right) - 3.$$
% Using again \eqref{binet} -- in particular the fact $\tau^{h+3}\geq (T_h-1/2)/C_1$ -- we deduce from \eqref{g} the general gap estimate
\begin{align*}\gamma_{3,N}&\geq
\frac{1}{C_1\tau^3} -\frac{1}{N}\left( \frac{1}{2C_1\tau^3}\log_\tau \left( \frac{N +\frac{1}{2}}{C_1} \right)+2\tau^2-2-\frac{1}{C_1\tau^3})\right).
\end{align*}

In the particular case $N=2$, we note that the subwords of length 2 in $\mathbf v_3$ are the following $\{11,12,13,21,31\}$, see \cite{seminalrauzy}. We then have
$$\lambda_{k+2}(\tau)-\lambda_k(\tau)\geq d_3(1)+d_3(3)=\tau^2-\tau.$$
\end{proof}


\section{Some remarks on the density of $\Lambda(q_m)$}\label{sec4}

The gap frequency of the spectrum $\Lambda(q)$ is deeply related to its \emph{upper density} \cite{beu}

\begin{equation*}
D^+(\Lambda(q)) := \lim_{R\to +\infty} \frac{n^+(R)}{R}
\end{equation*}
where $n^+(R)$ denotes the largest number of elements of $\Lambda(q)$ contained in an interval of length $R$, see \cite{komlorbook} for the existence of the limit. In \cite{FW02} Feng and Weng proved that every Pisot spectrum is associated with a substitution. When the substitution has the further property of being \emph{primitive} -- namely the power $M^k$ of its incidence matrix $M$ has positive entries for all sufficiently large $k$--  the associated symbolic dyanamics is strictly ergodic, from which follows that the gaps admit \emph{uniform frequencies}. Therefore the average gap length is independent of the position of the observation window and the upper density is given by
\begin{equation}\label{density}
D^+(\Lambda(q)) =\frac{1}{\sum_{k=1}^{N(q)} f_k(q)g_k(q)}.
\end{equation}
where $(f_k(q),g_k(q))$ are the finite, say $N(q)$, (uniform) frequency-gap couples associated to $\Lambda(q)$. In general, it is not known whether every Pisot number yield a uniform frequency of the gaps in the its spectrum, but several classes of Pisot numbers, including the $m$-bonacci numbers, with this property were provided.  
% \begin{remark}
% The formula \eqref{lambdaest} can be proved by the following arguments. 
% \begin{enumerate}
% \item $m$-bonacci words are balanced, so there exists a constant $b_m$ such that 
% $||w|_j-|u|_j|\leq b_m$
% for every couple of subwords of $\mathbf v$ of the same length and for all digit $j=1,\dots,m$.  Assuming that $u$ is a prefix of $\mathbf v_m$, dividing by $L:=|u|=|w|$ we obtain 
% $\frac{|w|_j}{L}=f_j(q_m)+o(L)$
% as $L\to \infty$.
% \item 
% \end{enumerate}
% \end{remark}
For instance, in the golden mean case $m=2$, such couples are $(\frac{1}{\varphi},1), (\frac{1}{\varphi^2},\varphi-1)$, $N(q_2)=2$ and the upper density is $D^+(\Lambda(\varphi))=\frac{1}{\frac{1}{\varphi}+\frac{1}{\varphi^3}}=\frac{\varphi^2}{\sqrt{5}}$ -- this can also be deduced by the fact that every sufficienly large $\lambda$ belongs to $\Lambda(\varphi)$ if and only if $ \lambda/\varphi$ belongs to the celebrated Fibonacci quasicrystal \cite{H04,MPP15}, whose density is $\frac{\varphi}{\sqrt{5}}$. 

 In \cite{GH06} the following result is proved. Let $q$ be the Pisot root of 
$$x^m-ax^{m-1}-\cdots ax -b$$
where $a\geq b\geq 1$. Let $r=\lfloor q\rfloor=a$. Then the gaps in $\Lambda^{r}(q)$ are 
$$ 1,\quad q-a, \quad q^2-a q-a,\quad\dots \quad  q^{m-1}-a q^{m-2}-\cdots a q-a$$ and they respectively 
occur with uniform frequency 
$$\frac{q^{m-1}}{Q},\,\frac{q^{m-2}}{Q},\dots, \frac{1}{Q}$$
where $Q:=q^{m-1}+q^{m-2}+\cdots+q+1=\frac{q^m-1}{q-1}$. This allows to compute the upper density 
\begin{align*}D^+(\Lambda^r(q))&=
% \frac{Q}
% {q^{m-1}+q^{m-2}(q-a)+q^{m-3}(q^2-aq-a)+\cdots+1\cdot(q^{m-1}-a q^{m-2}-\cdots a q-a)}\\
\frac{Q}
{mq^{m-1}-a\sum_{k=0}^{m-2}(k+1)q^{k}}\\
&=\frac{Q}{mq^{m-1}-a\frac{(m-1)q^{m}-mq^{m-1}+1}{(q-1)^2}}\\
&=\frac{(q^m-1)(q-1)}{mq^{m+1}-(2m+am-a)q^{m}+m(a+1)q^{m-1}-a}.
\end{align*}
In the case of $m$-bonacci Pisot numbers, the above formula yields 
$$D^+(\Lambda(q_m))= \frac{1}{\frac{1}{q_m}+\sum_{k=2}^m \frac{q_m^{k-1}-q_m^{k-2}-\cdots-q_m-1}{q_m^k}}=\frac{(q^m_m-1)(q_m-1)}{mq_m^{m+1}-(3m-1)q_m^{m}+2mq_m^{m-1}-1}.$$
% This can be obtained also by combining the Bugeaud's result \cite{bug02} and the formula \eqref{density}. 


The density provides an average distribution of the gaps. 
\medskip
 By \cite{bkl98}, see also \cite{komlorbook}, since $\Lambda(q_m)$ has finite upper density, then there exists $\gamma_{N,m}$ satisfying  \eqref{gamma} and 
\begin{equation}\label{denlim}\lim_{N\to \infty}\gamma_{N,m}=\frac{1}{D^+(\Lambda(q_m))}=\frac{1}{q_m}+\sum_{k=2}^m \frac{q_m^{k-1}-q_m^{k-2}-\cdots-q_m-1}{q_m^k}.\end{equation}
A trivial choice for $\gamma_{N,m}$ is the the smallest gap in $\Lambda(q_m)$, that is $\frac{1}{q_m^{m-1}}$, because 
$$\lambda_{k+N}(q_m)-\lambda_{k}(q_m)=\sum_{h=0}^{N-1}\lambda_{k+h+1}(q_m)-\lambda_{k+h}(q_m)\geq \frac{N}{q_m^{m-1}}.$$
Such an estimate is far from the “ideal” choice of $\gamma_{N,m}$ suggested by \eqref{denlim}, since it only reflects the minimal gap and ignores the finer distribution of gaps in the spectrum. An explicit estimate consistent with \eqref{denlim} for $m=2$ is for instance the estimate provided in \eqref{lambdaest}. 

\section{Conclusions}\label{sec5}
In this paper we established  explicit gap estimates for the spectrum of the  $m$-bonacci Pisot numbers. Our methodology allows for a finer investigation of the structure of the spectra and it relies on the combinatorical properties of the underlying $m$-bonacci words, particularly the existence of forbidden subwords, the balance property (factors of the same lengths have bounded variations in their digit distributions), and the digital representations of integers via $m$-bonacci number systems.  We then applied this general framework to derive further estimates for the Fibonacci and Tribonacci cases. Some considerations on the density of the spectra conclude our investigation. 


\bibliographystyle{abbrv}
\bibliography{rauzy}
\end{document}
